\subsection{Preprocessing and Feature Decisions}
\label{sec:preprocessing}

Both datasets are prepared the same way: the inputs are standardized on the training folds, the charge target is $\log_{10}(\Dch/M)$, the mass target is $M$ in solar masses, and whole curves are split into a frozen test set and GroupKFold folds. \Cref{tab:decisions-ns,tab:decisions-bh} list each decision with its grounds from Sections~\ref{sec:eda-ns} and~\ref{sec:eda-bh}; the three choices a table row cannot carry are argued first.

\paragraph{The charge keeps every value: no floor.}
The charge target is $\log_{10}(\Dch/M)$. Every prepared charge is positive -- the smallest neutron-star charge is $\nsEdaDMin$, the smallest black-hole charge $\bhEdaDMin$ -- so the logarithm needs no floor. The smallest neutron-star charges are not a numerical zero: they open their curves at low central density and grow smoothly and monotonically along them, so a floor would flatten real data into a constant and put a kink into the target. On the logarithmic scale every charge carries the same relative weight across the $\nsEdaDOrders$ orders of magnitude $\Dch$ spans, which is what a relative-error metric asks for (Section~\ref{sec:hypotheses}).

\paragraph{Whole curves are split: a frozen test set, then GroupKFold folds.}
Rows of one curve are each other's nearest neighbours ($\nsEdaSameCurvePercent\%$ of the neutron-star rows, $\bhEdaSameCurvePercent\%$ of the black-hole rows), so a split of rows would put near-copies of every test row into training and overstate accuracy (Sections~\ref{sec:eda-ns} and~\ref{sec:eda-bh}). Every split therefore assigns whole curves. First, $\prepNsTestPercent\%$ of the curves are frozen as the test set, chosen with the fixed seed $\prepSeed$: $\prepNsTestCurves$ of the $\prepNsCurves$ neutron-star curves and $\prepBhTestCurves$ of the $\prepBhCurves$ black-hole curves. The test set is used once, for the final evaluation. The remaining curves are shuffled with the same seed into $\prepNsFolds$ GroupKFold folds for model selection, $\prepNsFoldMin$ to $\prepNsFoldMax$ neutron-star curves and $\prepBhFoldMin$ to $\prepBhFoldMax$ black-hole curves per fold. The standardization statistics are fitted on the training folds only, so no information from a held-out curve reaches the model.

\paragraph{Extrapolation is measured in a separate ablation, not claimed.}
The final model is trained on all curves outside the test set, so it covers the whole sampled domain, and its validity is claimed only there. How it degrades beyond the domain is measured apart: an ablation trains without the curves on the rim of the domain ($\prepNsAblationCurves$ neutron-star curves on the edges of the $(\beta, \lambda)$ grid and along its empty corner; the $\prepBhAblationCurves$ outermost black-hole curves) and tests on them. The data already indicate the outcome: for neutron stars, predicting the rim from the interior costs about six times the error of predicting an interior curve ($\nsEdaSplitRimDM$ against $\nsEdaSplitCurvesDM$), while for black holes the outermost curves cost no more than any other ($\bhEdaSplitOuterCurvesLD$ against $\bhEdaSplitCurvesLD$).

% The decisions tables' layout is paper.toml [tables] decisions (W-124): wide or narrow.
\input{43_preprocessing_tab_decisions_ns_\prepDecisionsLayout}
\FloatBarrier

The black-hole decisions (\cref{tab:decisions-bh}) differ from the neutron-star ones in three rows: the horizon radius enters linearly, because it is uniform and spans only $\bhEdaRhOrders$ orders of magnitude; the charge target keeps the neutron-star definition although dividing by the mass gains nothing here, so that both datasets share one target; and the domain is bounded by the existence edge instead of a triangle.

\input{43_preprocessing_tab_decisions_bh_\prepDecisionsLayout}
\FloatBarrier
